\chapter{Словарь основных объектов}
\label{ch:definitions}

Эта глава вводит все базовые понятия \emph{до} того, как они используются в
механизмах. Определения намеренно сжаты — детали в соответствующих главах; здесь
задача, чтобы читатель (в том числе новый в проекте) владел словарём.

\section{Субстрат}

\begin{definition}[Граф памяти]
\emph{Граф памяти} $G=(V,E)$ — направленный граф, где $V$ — узлы-воспоминания, а
$E$ — типизированные знаковые взвешенные рёбра-отношения. Граф \emph{темпоральный}
(узлы несут время/слой) и \emph{динамический} (операторы меняют и узлы, и рёбра).
\end{definition}

\begin{definition}[Узел памяти]
\emph{Узел} $v\in V$ — запись с содержанием $\mathbf{x}_v$ и полями состояния:
вид \texttt{node\_kind} (observation/claim/hypothesis/preference/reaction\_pattern/
relationship\_state/\dots; задаёт смысл уверенности, гл.~\ref{ch:substrate}),
уверенность $\conf(v)\in[-1,1]$, валентность (аффективная окраска), временной слой,
провенанс, важность \texttt{importance}, флаги
\texttt{need\_to\_discover}/\texttt{need\_to\_revise}, транзиентная активация $a_v$.
Узел одновременно — данные (что помнится), убеждение (насколько истинно) и след
(насколько активно и свежо).
\end{definition}

\begin{definition}[Ребро и два его канала]
\emph{Ребро} $(u\to v)$ направлено и несёт \emph{тип} отношения (\textsc{isA},
\textsc{supports}, \textsc{contradicts}, \textsc{associates}, \dots). На нём живут
\textbf{два независимых числовых канала}, которые читают разные операторы:
\begin{itemize}
\item \textbf{вес перехода} $w_{uv}\ge 0$ — для обхода (walker); нормирован по узлу,
$\sum_{v} P(u\to v)=1$ (это и есть PageRank-подобная часть);
\item \textbf{знак} $s_{uv}\in\{+1,-1\}$ — полярность отношения (подтверждает /
опровергает); используется \emph{только} в пропагации уверенности
(глава~\ref{ch:propagation}), чтобы уверенность могла не только расти, но и убывать.
\end{itemize}
Знак и вес не конфликтуют: первый — про \emph{убеждение}, второй — про \emph{обход}.
\end{definition}

\begin{definition}[Уверенность]
$\conf(v)\in[-1,1]$ — \emph{знаковая} степень убеждённости: $+1$ — уверена, что
\emph{да}; $-1$ — уверена, что \emph{нет}; $0$ — не знает. Магнитуда $|\conf(v)|$ —
собственно уверенность, знак — направление убеждения. Эквивалентно вероятности
через $\conf=2P-1$ (где $P$ — вероятность истинности, $P=0{,}5\Leftrightarrow
\conf=0$), но эргономичнее: зона незнания — это $|\conf(v)|<h$ (порог $h\approx
0{,}3$), и такие узлы — кандидаты на \texttt{need\_to\_discover} или удаление.
Меняется пропагацией (глава~\ref{ch:propagation}) и ревизией
(глава~\ref{ch:epistemics}). Принцип: вывод не получает $|\conf|=1$.
\end{definition}

\begin{definition}[Активация]
\emph{Активация} $a_v(t)\ge 0$ — транзиентное скалярное поле «насколько узел
подсвечен \emph{сейчас}». Обновляется распространением активации
\[
a(t{+}1)=\lambda\,\tilde P^{\top} a(t)+s(t),\qquad \lambda\in(0,1),
\]
где $\tilde P$ — нормированная матрица переходов, $s(t)$ — вброс (визиты роя,
инъекция запроса), $\lambda$ — распад. Эквивалентно — ожидаемая частота посещений
узла роем в коротком окне. \textbf{Не путать:} активация — «здесь и сейчас»,
уверенность — медленное убеждение, стационарное $\pi$ — долгосрочная частота
(глава~\ref{ch:topology}).
\end{definition}

\begin{definition}[Провенанс (обоснование)]
\emph{Провенанс} узла — \emph{множество идентификаторов узлов-посылок}
$\{v_1,\dots,v_n\}$, из которых узел непосредственно выведен (для наблюдаемого факта
— идентификатор внешнего источника). Это \emph{justification set} в смысле
поддержки истинности (глава~\ref{ch:epistemics}), а \textbf{не} полный трек обхода:
трек шумен, доступен отдельно и обоснованием не является. Провенанс делает возможной
ретракцию: опровергнув посылку, находим и понижаем её следствия.
\end{definition}

\begin{definition}[Валентность и настроение]
\emph{Валентность} узла — его \emph{аффективная} окраска (пара
валентность–возбуждение), \textbf{ортогональная уверенности}: уверенность — насколько
узел \emph{истинен}, валентность — как он \emph{ощущается}. Пример: «точно знаю, что
это ужасное воспоминание» — $\conf$ высокая, валентность резко негативная.
Валентность — канал сцепки с настроением (mood-congruent recall,
глава~\ref{ch:mood}); уверенность — канал сцепки с пропагацией. \emph{Настроение}
$\mood$ — отдельный низкоразмерный вектор состояния Рин, сцепленный с графом
двусторонне, но не хранящийся в узлах.
\end{definition}

\begin{definition}[Время: системное и содержательное]
На узле хранится \emph{минимум системного времени}: тег \texttt{layer} (сессия
\dots жизнь) и одна метка (кодирование / последний доступ) — этого достаточно
операторам забывания и консолидации. \emph{Содержательное} время («когда это
произошло») моделируется как обычные узлы-факты — оно реконструктивно и может быть
неверным, как у человека. Интервал валидности (start/end) вешается \emph{только} на
факты, которым он нужен (квалификаторы Wikidata-схемы, глава~\ref{ch:substrate}), а
не на каждый узел.
\end{definition}

\section{Обход}

\begin{definition}[Случайный блуждающий агент, Random Walker]
\emph{Random Walker} (RW) — агент, который стартует в некотором узле и шаг за шагом
переходит к соседям по \emph{вероятностному} правилу
$P(A\to B)\propto w_{AB}\cdot\text{rel}(B)\cdot\text{mc}(B)\cdot\text{unc}(B)^{\beta_u}$
(глава~\ref{ch:walker}), оставляя за собой \emph{след} — последовательность
посещённых узлов. По ходу он может запускать операции (извлечение, вывод,
инспекцию).
\end{definition}

\noindent Функционал одного RW: (i)~перейти по вероятностному правилу; (ii)~с
вероятностью $\alpha$ \emph{телепортироваться} (рестарт в случайный узел);
(iii)~повысить активацию посещённого узла; (iv)~по условию запустить операцию
(retrieval / fusion / инспекцию need\_to\_*); (v)~записать след.

\begin{definition}[Рой блуждающих агентов, RW Swarm]
\emph{Swarm} — множество RW, запущенных параллельно с общей политикой. Рой
аппроксимирует распространение активации (spreading activation) и — на большой
статистике следов — эмпирическое распределение маршрутов по графу
(глава~\ref{ch:experiment}).
\end{definition}

\begin{definition}[Три политики обхода]
Один и тот же движок роя работает в трёх режимах, задаваемых параметрами
(смещение перехода, телепорт $\alpha$, температура): \emph{retrieval} (под
запрос), \emph{внутренняя жизнь/любопытство} (к неопределённости) и \emph{сон}
(стохастическое покрытие с телепортом). См.\ главу~\ref{ch:walker}.
\end{definition}

\section{Операции над памятью}

\begin{definition}[Retrieval (извлечение)]
\emph{Retrieval} — двухшаговая операция. \textbf{Шаг~1 (якоря):} быстрым индексом
(векторный/эмбеддинг-поиск, как в A-MEM/Zep) находим узлы, релевантные запросу; RW
для \emph{поиска} входа не годится (медленен, диффузен). \textbf{Шаг~2
(расширение):} от якорей запускаем ограниченный mood-смещённый рой RW, собирающий
активированную \emph{окрестность} — ассоциации, которых векторный поиск не даёт.
Итог зависит от настроения (mood-congruent recall). RW, таким образом, нужен и днём
(расширение), а сон — тот же RW плюс телепорт и без границ.
\end{definition}

\begin{definition}[Fusion (слияние, вывод)]
\emph{Fusion} — порождение \emph{нового узла} из совместно активированных: система
замечает согласованную конфигурацию посылок и создаёт узел-\emph{гипотезу} с низкой
уверенностью и провенансом-обоснованием. Fusion — механизм рождения выводов и
ассоциаций (глава~\ref{ch:fusion}).
\end{definition}

\begin{definition}[Консолидация]
\emph{Консолидация} (в широком смысле) — \emph{процесс} перехода от множества частных
следов к устойчивым обобщениям, а не единичный оператор. Он отрабатывает несколько
подоператоров: fusion\,/\,вывод, summary-overlay, dedup\,/\,merge и — как \emph{поздний}
подоператор — prune (destructive replacement с огрублением, graph coarsening).
Большинство ранних консолидаций ничего не разрушают; разбор процесса и его
предохранителей — в главе~\ref{ch:timescales}.
\end{definition}

\begin{definition}[Забывание (decay)]
\emph{Забывание} — ослабление неиспользуемого; отрицательная сторона консолидации.
Моделируется базовой активацией (частота $+$ свежесть, глава~\ref{ch:timescales}).
\end{definition}

\begin{definition}[Ночная работа (сон)]
\emph{Сон} — фоновый режим, когда телепортирующий рой случайно инспектирует память:
консолидирует, проверяет противоречия (скраб), рождает неожиданные fusion-ы,
подтягивает внешнее по флагам \texttt{need\_to\_discover} (глава~\ref{ch:sleep}).
\end{definition}

\begin{definition}[Значимость (salience) и оценивание]
\emph{Значимость} события — мера того, насколько оно важно \emph{для самой Рин}
(теория оценивания). Дешёвый прокси — величина изменения настроения $|\Delta\mood|$;
глубокий сигнал — ошибка предсказания. Значимость гейтит консолидацию
(глава~\ref{ch:personality}).
\end{definition}

\section{Выражение}

\begin{definition}[Личность]
\emph{Личность} — не таблица весов в графе, а распределённый процесс: мотивация
$+$ настроение $+$ значимость всплывшего, входящий как \emph{смещение во время
обхода}, а на выходе — как стиль (глава~\ref{ch:personality}).
\end{definition}

\begin{definition}[Rinness / голос]
\emph{Голос} — языковая модель, превращающая внутренний контекст (всплывшие
ингредиенты, настроение, выбранный тон) в реплику в узнаваемом стиле Рин. Часть
характера живёт здесь; факты — в графе (глава~\ref{ch:voice}).
\end{definition}

\begin{definition}[Эмбриональный граф]
\emph{Эмбриональный граф} — начальное наполнение памяти из готовых баз здравого
смысла (ConceptNet $+$ ATOMIC), нейтральное по характеру; на нём Рин «взращивает»
собственную память фьюжнами и консолидацией (глава~\ref{ch:embryonic}).
\end{definition}
